\documentclass[12pt,a4paper]{article}

\usepackage[utf8]{inputenc}
\usepackage[T1]{fontenc}
\usepackage[margin=1in]{geometry}
\usepackage{array}
\usepackage{booktabs}
\usepackage{hyperref}
\usepackage{longtable}
\usepackage{tabularx}
\usepackage{xcolor}
\usepackage{listings}
\usepackage{float}
\usepackage{graphicx}

\hypersetup{
  colorlinks=true,
  linkcolor=blue,
  urlcolor=blue,
  citecolor=blue
}

\lstdefinestyle{fighterstyle}{
  basicstyle=\ttfamily\small,
  breaklines=true,
  columns=fullflexible,
  frame=single,
  backgroundcolor=\color{gray!5}
}

\title{Final Project Design Document}
\author{}
\date{April 28, 2026}

\begin{document}

\maketitle

This document matches the current repository state as of April 28, 2026.
Compared with the earlier draft, the most important change is that the
project no longer treats FPGA VGA as only future work. The current hardware
build includes both a custom FPGA VGA renderer and a custom WM8731
audio peripheral, each exposed to HPS Linux through the lightweight
HPS-to-FPGA bridge and accessed from user space through \texttt{/dev/mem}.

\section{Introduction and Overview}

\subsection{Project Summary}

This project implements a two-player fighting game prototype on the
\texttt{DE1-SoC} platform. The current system combines HPS-side game logic
with FPGA-side video and audio output:

\begin{itemize}
  \item The \texttt{HPS} side in C handles USB keyboard input, game-state
  updates, collision and damage logic, animation selection, software scene
  composition, audio command scheduling, and MMIO access.
  \item The primary video path is now the FPGA VGA renderer
  \texttt{fighter\_vga\_renderer}, not the Linux \texttt{/dev/fb0}
  framebuffer. HPS writes \texttt{320x240} RGB565 frame data through
  Avalon-MM. The FPGA owns VGA timing, display buffering, RGB output, and
  vblank-safe double-buffer swaps.
  \item The video code still includes fallbacks to Linux framebuffer devices
  such as \texttt{/dev/fb0} and to console output when the FPGA VGA path is
  unavailable.
  \item The FPGA side implements a \texttt{WM8731} audio peripheral,
  \texttt{fighter\_audio\_wm8731}, with FIFO-backed sample writes, I2C codec
  initialization, and serial DAC output.
  \item The software still contains \texttt{fighter\_mmio\_encode()}, which
  packs game state into 22 words. In the current hardware design, however, the
  implemented VGA IP accepts MMIO frame data rather than directly consuming
  those 22 game-state words.
\end{itemize}

\subsection{Relationship Between Current Implementation and Target Architecture}

\begin{longtable}{>{\raggedright\arraybackslash}p{0.28\textwidth} >{\raggedright\arraybackslash}p{0.18\textwidth} >{\raggedright\arraybackslash}p{0.42\textwidth}}
\toprule
Layer & Current Status & Notes \\
\midrule
\endhead
Dual USB keyboard input & Implemented & \texttt{libusb} polling for up to two HID boot keyboards when \texttt{libusb-1.0} is available \\
Fighting game state machine & Implemented & Menu, match, game over, KO, timeout, exit, blocking, projectiles, hit and block resolution \\
Animation and assets & Implemented & Ryu and Ken PPM sprite sets, background image, menu images, and fireball projectiles \\
FPGA VGA renderer & Implemented & \texttt{fighter\_vga\_renderer} at default physical address \texttt{0xFF240000}; this is the main display path \\
Linux framebuffer / console video fallback & Implemented & Used only if VGA MMIO probing fails or if console mode is requested \\
FPGA audio peripheral & Implemented & \texttt{fighter\_audio\_wm8731} at default physical address \texttt{0xFF200000} \\
Game-state MMIO encoder & Implemented in software & Exports 22 32-bit words, but current VGA hardware uses a framebuffer protocol instead \\
Linux kernel driver & Not implemented & Current board access is through user-space \texttt{/dev/mem} and optional \texttt{/dev/fb*} \\
\bottomrule
\end{longtable}

\subsection{Assumptions and Constraints}

\begin{itemize}
  \item The target board is \texttt{DE1-SoC}.
  \item The system assumes up to two USB keyboards by default.
  \item The default input mapping is \texttt{W/A/S/D/J/K/L}.
  \item The main logic runs at approximately 60 FPS.
  \item The game coordinate system is \texttt{640x480}. The FPGA framebuffer
  is \texttt{320x240} and is doubled to \texttt{640x480} by hardware.
  \item The current design uses direct user-space physical memory mapping
  instead of a Linux kernel driver.
\end{itemize}

\section{Top-Level System Block Diagram}

\subsection{Actual Current Data Path}

\begin{lstlisting}[style=fighterstyle]
USB Keyboard 1/2
      |
      v
usb_hid_keyboard.c  (libusb, up to 2 devices)
      |
      v
fighter_input.c
      |
      v
fighter_game.c  -----> fighter_animation.c
      |                         |
      | game state              | current sprites
      v                         v
fighter_renderer.c  ----------------------------+
      |                                         |
      | RGB565 frame writes through MMIO        |
      v                                         |
/dev/mem + lw HPS-to-FPGA bridge                |
      |                                         |
      v                                         |
fighter_vga_renderer (FPGA VGA renderer)        |
      |                                         |
      | VGA timing / buffer swap / RGB output   |
      v                                         |
VGA_R/G/B, HS, VS, CLK                          |

fighter_game.c
      |
      | audio commands
      v
fighter_audio.c
      |
      +--> /dev/mem + lw bridge --> fighter_audio_wm8731 --> WM8731
      |
      +--> command backend fallback when MMIO audio is unavailable
\end{lstlisting}

\subsection{Platform Address Map}

\texttt{hw/soc\_system.qsys} connects both project-specific peripherals to
the HPS lightweight bridge. With the standard lightweight bridge physical
base of \texttt{0xFF200000}, the current software defaults are:

\begin{longtable}{>{\raggedright\arraybackslash}p{0.22\textwidth} >{\raggedright\arraybackslash}p{0.20\textwidth} >{\raggedright\arraybackslash}p{0.22\textwidth} >{\raggedright\arraybackslash}p{0.24\textwidth}}
\toprule
IP Instance & Qsys Base & Linux Physical Address & Override Variable \\
\midrule
\endhead
\texttt{fighter\_audio\_0} & \texttt{0x0000} & \texttt{0xFF200000} & \texttt{FIGHTER\_AUDIO\_MMIO\_ADDR} \\
\texttt{fighter\_vga\_0} & \texttt{0x40000} & \texttt{0xFF240000} & \texttt{FIGHTER\_VGA\_MMIO\_ADDR} \\
\bottomrule
\end{longtable}

Both VGA and audio software also access the HPS bridge reset register at
\texttt{0xFFD0501C}. The software clears bits \texttt{[1:0]} to enable the
HPS-FPGA bridge before probing the custom IP.

\subsection{Image and Audio Assets}

Game images and sounds are stored as normal files in the HPS Linux filesystem.
They are not compiled into the FPGA bitstream. The main VGA path is still the
custom FPGA VGA renderer: HPS writes RGB565 frame data through MMIO, and the
FPGA handles VGA timing, buffer swapping, and RGB output. For audio, HPS writes
PCM samples into audio MMIO FIFOs, and the FPGA drives the WM8731.

\begin{table}[H]
\footnotesize
\setlength{\tabcolsep}{4pt}
\begin{tabularx}{\textwidth}{>{\raggedright\arraybackslash}p{0.20\textwidth} >{\raggedright\arraybackslash}X >{\raggedright\arraybackslash}p{0.22\textwidth}}
\toprule
Asset Type & Repository Path & Runtime Loader \\
\midrule
Ryu sprites & \path|game_assets/sprites/RyuPPM/*/*.ppm| & \path|fighter_animation.c| \\
Ken sprites & \path|game_assets/sprites/KenPPM/*/*.ppm| & \path|fighter_animation.c| \\
Fireball sprites & \path|game_assets/sprites/*PPM/fireball/*.ppm| or \path|attack_fireball/*.ppm| & \path|fighter_animation.c| \\
Background & \path|game_assets/background/background.ppm| & \path|fighter_renderer.c| \\
Menu frames & \path|game_assets/ui/menu/menu_frame_0.ppm| and \path|menu_frame_1.ppm| & \path|fighter_renderer.c| \\
Menu BGM & \path|game_assets/sound effects/Title.wav| & \path|fighter_audio.c| \\
Confirm sound & \path|game_assets/sound effects/Credit.wav| & \path|fighter_audio.c| \\
Game-over sound & \path|game_assets/sound effects/Game Over.wav| & \path|fighter_audio.c| \\
\bottomrule
\end{tabularx}
\end{table}

The full loader source paths are \texttt{sw/fighter\_animation.c},
\texttt{sw/render\_if/fighter\_renderer.c}, and
\texttt{sw/audio/fighter\_audio.c}.

The runtime asset search rules are:

\begin{itemize}
  \item If \texttt{FIGHTER\_ASSET\_ROOT} is set, animation code uses
  \texttt{\$FIGHTER\_ASSET\_ROOT/sprites/RyuPPM} and
  \texttt{\$FIGHTER\_ASSET\_ROOT/sprites/KenPPM}. The renderer uses the same
  root for background and menu images.
  \item If \texttt{FIGHTER\_ASSET\_ROOT} is not set, animation first tries
  \texttt{/root/game\_assets/sprites/RyuPPM} and
  \texttt{/root/game\_assets/sprites/KenPPM}, then falls back to the repository
  relative paths \texttt{../game\_assets/sprites/RyuPPM} and
  \texttt{../game\_assets/sprites/KenPPM}.
  \item Background/menu loading similarly tries \texttt{FIGHTER\_ASSET\_ROOT},
  \texttt{/root/game\_assets}, and \texttt{../game\_assets}.
  \item The current WAV path function returns repository-relative paths under
  \texttt{../game\_assets/sound effects}. The MMIO audio backend opens them
  with \texttt{fopen()}, parses RIFF/WAVE chunks, accepts PCM16 mono or stereo,
  and resamples to the software target rate of \texttt{48000} Hz.
\end{itemize}

PPM image loading supports both \texttt{P6} binary and \texttt{P3} ASCII PPM.
The software stores decoded pixels as 24-bit RGB in heap memory. When drawing
to the FPGA VGA backend, the renderer packs the final composed frame into
RGB565 words and writes those words through MMIO.

\section{Custom Peripheral Register Maps}

\subsection{Implemented VGA MMIO Contract}

The current video hardware source is \texttt{hw/fighter\_vga\_renderer.sv},
packaged by \texttt{hw/fighter\_vga\_hw.tcl}. Its Avalon-MM interface uses
32-bit data and word-based addresses. Software therefore accesses
\texttt{regs[n]} for Avalon word offset \texttt{n}.

\begin{longtable}{>{\raggedright\arraybackslash}p{0.15\textwidth} >{\raggedright\arraybackslash}p{0.12\textwidth} >{\raggedright\arraybackslash}p{0.14\textwidth} >{\raggedright\arraybackslash}p{0.47\textwidth}}
\toprule
Avalon Signal & Width & Direction & Meaning \\
\midrule
\endhead
\texttt{avs\_chipselect} & 1 & HPS to FPGA & Selects this slave for the current transaction \\
\texttt{avs\_read} & 1 & HPS to FPGA & Indicates a register read transaction \\
\texttt{avs\_write} & 1 & HPS to FPGA & Indicates a register or framebuffer write transaction \\
\texttt{avs\_address} & 16 & HPS to FPGA & 32-bit word offset; \texttt{0}, \texttt{1}, \texttt{2}, etc., not byte address \\
\texttt{avs\_writedata} & 32 & HPS to FPGA & Control word or two packed RGB565 pixels \\
\texttt{avs\_readdata} & 32 & FPGA to HPS & Control/status, geometry, ident, or zero for unmapped addresses \\
\bottomrule
\end{longtable}

\begin{longtable}{>{\raggedright\arraybackslash}p{0.14\textwidth} >{\raggedright\arraybackslash}p{0.22\textwidth} >{\raggedright\arraybackslash}p{0.10\textwidth} >{\raggedright\arraybackslash}p{0.40\textwidth}}
\toprule
Word Offset & Register & R/W & Description \\
\midrule
\endhead
\texttt{0} & \texttt{REG\_CONTROL} & RW & Status and frame-swap control \\
\texttt{1} & \texttt{REG\_WIDTH} & R & Returns framebuffer width \texttt{320} \\
\texttt{2} & \texttt{REG\_HEIGHT} & R & Returns framebuffer height \texttt{240} \\
\texttt{3} & \texttt{REG\_STRIDE} & R & Returns row stride \texttt{640} bytes \\
\texttt{31} & \texttt{REG\_IDENT} & R & Returns \texttt{0x56504741}, ASCII \texttt{"VPGA"} \\
\texttt{1024..39423} & frame input window & W & RGB565 frame data for the FPGA VGA renderer, two pixels per 32-bit word \\
\bottomrule
\end{longtable}

\noindent \texttt{REG\_CONTROL} read bits:

\begin{itemize}
  \item \texttt{bit0}: IP present, fixed to \texttt{1}
  \item \texttt{bit1}: swap pending
  \item \texttt{bit8}: current display buffer
  \item \texttt{bit9}: current software write buffer
\end{itemize}

\noindent \texttt{REG\_CONTROL} write bits:

\begin{itemize}
  \item \texttt{bit1}: swap request. Software writes this after filling a full
  frame. Hardware performs the buffer swap at vblank and then clears
  \texttt{swap\_pending}.
\end{itemize}

\subsection{VGA Framebuffer Format}

\begin{itemize}
  \item Framebuffer width: \texttt{320} pixels
  \item Framebuffer height: \texttt{240} pixels
  \item Pixel format: RGB565
  \item Pixels per 32-bit word: \texttt{2}
  \item Framebuffer words: \texttt{38400}
  \item Framebuffer start word offset: \texttt{1024}
  \item Total software mapping span: \texttt{39424 * 4 = 157696} bytes
\end{itemize}

\begin{lstlisting}[style=fighterstyle]
word[15:0]  = first RGB565 pixel
word[31:16] = second RGB565 pixel
\end{lstlisting}

The FPGA output timing is standard \texttt{640x480}. Internally, the renderer
uses the 50 MHz board clock, toggles a pixel tick to generate approximately
25 MHz pixel timing, and maps each framebuffer source pixel to a \texttt{2x2}
area on the VGA output.

\subsection{FPGA VGA Renderer Transaction Protocol}

The VGA backend in \texttt{sw/render\_if/fighter\_renderer.c} uses the
following protocol:

\begin{enumerate}
  \item Open \texttt{/dev/mem}.
  \item Map the HPS bridge reset register at \texttt{0xFFD0501C} and clear
  bits \texttt{[1:0]}.
  \item Map the VGA MMIO span at \texttt{FIGHTER\_VGA\_MMIO\_ADDR}, default
  \texttt{0xFF240000}.
  \item Read \texttt{regs[31]} and require \texttt{0x56504741}
  (\texttt{"VPGA"}).
  \item Read \texttt{regs[1]}, \texttt{regs[2]}, and \texttt{regs[3]} and
  require \texttt{320}, \texttt{240}, and \texttt{640}.
  \item Before writing a new frame, poll \texttt{regs[0] bit1} until
  \texttt{swap\_pending} is clear.
  \item Pack every two RGB565 pixels into one 32-bit word, and write
  \texttt{38400} words starting at \texttt{regs[1024]}.
  \item Write \texttt{regs[0] = 0x00000002}. The FPGA VGA renderer flips the
  display buffer on the next vblank and continues generating VGA timing and
  RGB output.
\end{enumerate}

\subsection{VGA Pins}

\begin{longtable}{>{\raggedright\arraybackslash}p{0.18\textwidth} >{\raggedright\arraybackslash}p{0.12\textwidth} >{\raggedright\arraybackslash}p{0.12\textwidth} >{\raggedright\arraybackslash}p{0.44\textwidth}}
\toprule
Port & Width & Direction & Data \\
\midrule
\endhead
\texttt{vga\_r} & 8 & output & Red channel, expanded from RGB565 red bits \\
\texttt{vga\_g} & 8 & output & Green channel, expanded from RGB565 green bits \\
\texttt{vga\_b} & 8 & output & Blue channel, expanded from RGB565 blue bits \\
\texttt{vga\_hs} & 1 & output & Active-low horizontal sync \\
\texttt{vga\_vs} & 1 & output & Active-low vertical sync \\
\texttt{vga\_clk} & 1 & output & Pixel clock derived by toggling at 50 MHz input clock \\
\texttt{vga\_blank\_n} & 1 & output & High during visible region \\
\texttt{vga\_sync\_n} & 1 & output & Fixed low for the DE1-SoC VGA DAC \\
\bottomrule
\end{longtable}

\subsection{Implemented Audio MMIO Contract}

The current audio hardware source is \texttt{hw/fighter\_audio.sv}, module
\texttt{fighter\_audio\_wm8731}, packaged by
\texttt{hw/fighter\_audio\_hw.tcl}. The default physical address is
\texttt{0xFF200000}.

\begin{longtable}{>{\raggedright\arraybackslash}p{0.15\textwidth} >{\raggedright\arraybackslash}p{0.12\textwidth} >{\raggedright\arraybackslash}p{0.14\textwidth} >{\raggedright\arraybackslash}p{0.47\textwidth}}
\toprule
Avalon Signal & Width & Direction & Meaning \\
\midrule
\endhead
\texttt{avs\_chipselect} & 1 & HPS to FPGA & Selects this audio slave \\
\texttt{avs\_read} & 1 & HPS to FPGA & Reads \texttt{control}, \texttt{fifospace}, or zero from sample ports \\
\texttt{avs\_write} & 1 & HPS to FPGA & Writes clear commands or sample words \\
\texttt{avs\_address} & 2 & HPS to FPGA & 32-bit word offset \texttt{0..3} \\
\texttt{avs\_writedata} & 32 & HPS to FPGA & Control clear bits or \texttt{int16} sample in upper halfword \\
\texttt{avs\_readdata} & 32 & FPGA to HPS & Status and FIFO-space fields \\
\bottomrule
\end{longtable}

\begin{longtable}{>{\raggedright\arraybackslash}p{0.12\textwidth} >{\raggedright\arraybackslash}p{0.18\textwidth} >{\raggedright\arraybackslash}p{0.08\textwidth} >{\raggedright\arraybackslash}p{0.48\textwidth}}
\toprule
Word Offset & Name & R/W & Description \\
\midrule
\endhead
\texttt{0} & \texttt{control} & RW & Read FIFO and codec status; write clear commands \\
\texttt{1} & \texttt{fifospace} & R & \texttt{[31:24]} left write space, \texttt{[23:16]} right write space \\
\texttt{2} & \texttt{leftdata} & W & Left-channel sample word; software writes \texttt{int16 << 16} \\
\texttt{3} & \texttt{rightdata} & W & Right-channel sample word; software writes \texttt{int16 << 16} \\
\bottomrule
\end{longtable}

\noindent The current \texttt{control} read word uses:

\begin{itemize}
  \item \texttt{bit0}: \texttt{codec\_init\_done}
  \item \texttt{bit1}: \texttt{codec\_init\_error}
  \item \texttt{bit2}: \texttt{tx\_underflow\_seen}
  \item \texttt{bit3}: \texttt{write\_overflow\_seen}
  \item \texttt{bits[15:8]}: \texttt{right\_count}
  \item \texttt{bits[23:16]}: \texttt{left\_count}
\end{itemize}

\noindent The \texttt{control} write word uses:

\begin{itemize}
  \item \texttt{bit2}: clear read/playback-side FIFO state
  \item \texttt{bit3}: clear write-side FIFO state
\end{itemize}

\subsection{Audio Software-Hardware Transaction Protocol}

The MMIO audio backend in \texttt{sw/audio/fighter\_audio.c} operates as:

\begin{enumerate}
  \item Open \texttt{/dev/mem}.
  \item Map the HPS bridge reset register at \texttt{0xFFD0501C} and clear
  bits \texttt{[1:0]}.
  \item Map four 32-bit words at \texttt{FIGHTER\_AUDIO\_MMIO\_ADDR}, default
  \texttt{0xFF200000}.
  \item Write \texttt{control bit2 | bit3}, then write zero, to clear FIFO and
  sticky status.
  \item Read \texttt{fifospace}. A valid peripheral reports nonzero left and
  right write space.
  \item Load WAV assets from the filesystem, parse RIFF chunks, convert PCM16
  mono/stereo to the internal stereo stream, and resample toward
  \texttt{48000} Hz.
  \item A worker thread polls \texttt{fifospace}; when both channels have room,
  it writes the left sample to \texttt{regs[2]} and right sample to
  \texttt{regs[3]}.
\end{enumerate}

\subsection{Audio Pins}

\begin{longtable}{>{\raggedright\arraybackslash}p{0.20\textwidth} >{\raggedright\arraybackslash}p{0.10\textwidth} >{\raggedright\arraybackslash}p{0.12\textwidth} >{\raggedright\arraybackslash}p{0.44\textwidth}}
\toprule
Port & Width & Direction & Data \\
\midrule
\endhead
\texttt{aud\_xck} & 1 & output & WM8731 master clock, default \texttt{12.5 MHz} \\
\texttt{aud\_bclk} & 1 & output & Audio bit clock, default \texttt{3.125 MHz} \\
\texttt{aud\_daclrck} & 1 & output & DAC left/right word-select clock \\
\texttt{aud\_adclrck} & 1 & output & ADC left/right word-select clock, synchronized to DAC LRCK \\
\texttt{aud\_dacdat} & 1 & output & Left-justified serial DAC sample bitstream \\
\texttt{aud\_adcdat} & 1 & input & ADC serial data input; present but not consumed by current software \\
\texttt{fpga\_i2c\_sclk} & 1 & bidir & Open-drain style I2C SCL for codec initialization \\
\texttt{fpga\_i2c\_sdat} & 1 & bidir & Open-drain style I2C SDA for codec initialization and ACK sampling \\
\texttt{codec\_init\_done} & 1 & output & Codec initialization completed; also routed to \texttt{LEDR[0]} \\
\texttt{codec\_init\_error} & 1 & output & Codec initialization failed; also routed to \texttt{LEDR[1]} \\
\bottomrule
\end{longtable}

\subsection{Software Game-State MMIO Encoder}

\texttt{sw/render\_if/fighter\_mmio.c} still implements
\texttt{fighter\_mmio\_encode()}, and \texttt{sw/include/fighter\_mmio.h}
still defines a 22-word game-state register layout. This contract is useful
for tests and for a possible future hardware sprite engine, but it is not the
register map consumed by the current framebuffer VGA IP.

\begin{longtable}{>{\raggedright\arraybackslash}p{0.12\textwidth} >{\raggedright\arraybackslash}p{0.28\textwidth} >{\raggedright\arraybackslash}p{0.48\textwidth}}
\toprule
Offset & Register & Software Encoding Meaning \\
\midrule
\endhead
\texttt{0x00} & \texttt{GAME\_STATE} & \texttt{bit[1:0]} menu/playing/game-over, \texttt{bit8} menu frame, \texttt{bit9} game-over ready \\
\texttt{0x04} & \texttt{PLAYER1\_X} & Player 1 top-left x coordinate \\
\texttt{0x08} & \texttt{PLAYER1\_Y} & Player 1 top-left y coordinate \\
\texttt{0x0C} & \texttt{PLAYER1\_STATE} & Player 1 visual state \\
\texttt{0x10} & \texttt{PLAYER2\_X} & Player 2 top-left x coordinate \\
\texttt{0x14} & \texttt{PLAYER2\_Y} & Player 2 top-left y coordinate \\
\texttt{0x18} & \texttt{PLAYER2\_STATE} & Player 2 visual state \\
\texttt{0x1C} & \texttt{PLAYER1\_HP} & Player 1 HP \\
\texttt{0x20} & \texttt{PLAYER2\_HP} & Player 2 HP \\
\texttt{0x24} & \texttt{ROUND\_TIMER} & Remaining seconds \\
\texttt{0x28} & \texttt{WINNER} & Winner enum \\
\texttt{0x2C} & \texttt{PLAYER1\_FACING} & \texttt{1 = right}, \texttt{0 = left} \\
\texttt{0x30} & \texttt{PLAYER2\_FACING} & \texttt{1 = right}, \texttt{0 = left} \\
\texttt{0x34} & \texttt{DEBUG\_FLAGS} & Last attacks and attack phases \\
\texttt{0x38} & \texttt{PLAYER1\_ATTACK\_CMD} & Player 1 last attack command \\
\texttt{0x3C} & \texttt{PLAYER2\_ATTACK\_CMD} & Player 2 last attack command \\
\texttt{0x40} & \texttt{PLAYER1\_STATE\_FRAME} & Frames spent in current visual state \\
\texttt{0x44} & \texttt{PLAYER2\_STATE\_FRAME} & Frames spent in current visual state \\
\texttt{0x48} & \texttt{PLAYER1\_EVENT\_FLAGS} & Per-frame event pulses \\
\texttt{0x4C} & \texttt{PLAYER2\_EVENT\_FLAGS} & Per-frame event pulses \\
\texttt{0x50} & \texttt{PLAYER1\_COMBAT\_RESULT} & None, hit, blocked, trade, or whiff \\
\texttt{0x54} & \texttt{PLAYER2\_COMBAT\_RESULT} & None, hit, blocked, trade, or whiff \\
\bottomrule
\end{longtable}

\section{Detailed Hardware Design}

\subsection{Top-Level Hardware Structure}

The board-level top module is \texttt{hw/soc\_system\_top.sv}. The current
project-relevant connections are:

\begin{itemize}
  \item \texttt{soc\_system} provides the HPS skeleton and lightweight bridge.
  \item \texttt{fighter\_audio\_0} is exported through the audio conduit and
  wired to \texttt{AUD\_*} and \texttt{FPGA\_I2C\_*} pins.
  \item \texttt{audio\_init\_done} and \texttt{audio\_init\_error} are mapped
  to \texttt{LEDR[0]} and \texttt{LEDR[1]}.
  \item \texttt{fighter\_vga\_0} is exported through the VGA conduit and wired
  to \texttt{VGA\_R/G/B}, \texttt{VGA\_HS}, \texttt{VGA\_VS},
  \texttt{VGA\_CLK}, \texttt{VGA\_BLANK\_N}, and \texttt{VGA\_SYNC\_N}.
\end{itemize}

\subsection{Hardware Module Responsibilities}

\begin{table}[H]
\small
\setlength{\tabcolsep}{4pt}
\begin{tabularx}{\textwidth}{>{\raggedright\arraybackslash}p{0.31\textwidth} >{\raggedright\arraybackslash}X}
\toprule
Module & Responsibility \\
\midrule
\path|soc_system_top| & Board-level wrapper. It connects DE1-SoC pins to the generated Platform Designer system, routes VGA/audio conduits, maps codec status to LEDs, and ties unused board outputs to simple quiet values. \\
\path|soc_system| & Platform Designer system. It contains the HPS, lightweight HPS-to-FPGA AXI master, interconnect, reset controllers, \path|fighter_vga_0|, and \path|fighter_audio_0|. \\
\path|soc_system_mm_interconnect_0| & Avalon-MM interconnect generated by Platform Designer. It decodes lightweight bridge addresses and forwards reads/writes to the selected custom IP. \\
\path|fighter_vga_renderer| & Custom VGA IP. It exposes geometry/control registers, accepts framebuffer writes, tracks display/write buffers, generates VGA timing, and outputs 8-bit RGB plus sync signals. \\
\path|fighter_audio_wm8731| & Custom audio IP. It exposes a 4-word register map, queues left/right samples in FIFOs, initializes the WM8731 by I2C, derives audio clocks, and serializes samples to the DAC. \\
\bottomrule
\end{tabularx}
\end{table}

\subsection{Register Versus BRAM Implementation}

The design contains three different kinds of state, and they synthesize
differently:

\begin{itemize}
  \item Small control/status values such as \texttt{swap\_pending},
  \texttt{display\_buffer}, VGA counters, audio FIFO pointers, FIFO counts,
  and I2C state are normal flip-flop registers.
  \item Audio sample storage is inferred as BRAM. The Quartus map/fitter
  reports show \texttt{left\_fifo\_rtl\_0} and \texttt{right\_fifo\_rtl\_0} as
  \texttt{altsyncram} simple dual-port RAMs, each \texttt{128 x 32} bits. Each
  FIFO uses one M10K block, so audio uses two M10K blocks total.
  \item The source version of \path|fighter_vga_renderer.sv| describes two
  framebuffer RAM banks with \texttt{altsyncram}. The current Quartus fit
  report, \path|hw/output_files/soc_system.fit.rpt|, reports zero M10Ks under
  \path|fighter_vga_renderer| and only two M10Ks total, both under the audio
  FIFOs. Therefore the current compiled resource report confirms BRAM for
  audio FIFOs, not a full synthesized on-chip VGA framebuffer. The
  software/hardware register protocol is still documented above because it is
  the intended VGA IP contract in the source tree.
\end{itemize}

\subsection{Internal Structure of \texttt{fighter\_vga\_renderer}}

\begin{lstlisting}[style=fighterstyle]
Avalon-MM Slave
  |
  +--> control / geometry / ident registers
  +--> framebuffer write window
          |
          v
     dual M10K framebuffer banks
          |
          v
     vblank-safe buffer swap
          |
          v
     640x480 VGA timing and 2x scaling
          |
          v
     RGB565 to 8-bit VGA RGB output
\end{lstlisting}

The module uses two framebuffer banks. Software always writes the bank that is
not currently displayed. After a full frame is written, software writes
\texttt{CONTROL\_SWAP\_REQUEST}; hardware flips the displayed bank at the
start of vertical blanking.

\subsection{Internal Structure of \texttt{fighter\_audio\_wm8731}}

\begin{lstlisting}[style=fighterstyle]
Avalon-MM Slave
  |
  +--> control / status
  +--> left FIFO
  +--> right FIFO
          |
          v
     sample serializer ------> AUD_DACDAT
          |
          +--> clock dividers ---> AUD_XCK / AUD_BCLK / AUD_DACLRCK / AUD_ADCLRCK

I2C init FSM -----------------> FPGA_I2C_SCLK / FPGA_I2C_SDAT
                                |
                                v
                              WM8731
\end{lstlisting}

\subsection{Audio Peripheral Timing and Protocol}

\texttt{fighter\_audio\_wm8731} operates in a single 50 MHz clock domain.
With default parameters:

\begin{itemize}
  \item \texttt{aud\_xck = clk / 4 = 12.5 MHz}
  \item \texttt{aud\_bclk = clk / 16 = 3.125 MHz}
  \item \texttt{lrck = clk / 1024 = 48.828125 kHz}
\end{itemize}

The codec is configured for left-justified format, 16-bit samples, slave mode,
DAC playback enabled, and active output.

\section{Software Interfaces}

\subsection{Core Game Interface}

\texttt{sw/include/fighter\_game.h} defines the game data model.

Key enumerations include:

\begin{itemize}
  \item \texttt{fighter\_game\_state\_t}: \texttt{MENU / PLAYING / GAME\_OVER}
  \item \texttt{fighter\_character\_id\_t}: \texttt{RYU / KEN}
  \item \texttt{fighter\_visual\_state\_t}: includes \texttt{IDLE},
  \texttt{WALK}, \texttt{JUMP}, \texttt{CROUCH}, \texttt{GUARD},
  \texttt{ATTACK}, \texttt{HIT}, \texttt{BLOCK\_STUN}, \texttt{KO},
  \texttt{VICTORY}, and \texttt{CROUCH\_GUARD}
  \item \texttt{fighter\_attack\_phase\_t}: \texttt{NONE}, \texttt{STARTUP},
  \texttt{ACTIVE}, \texttt{HIT\_CONFIRM}, \texttt{BLOCK\_CONFIRM},
  \texttt{RECOVERY}
  \item \texttt{fighter\_winner\_t} and \texttt{fighter\_finish\_reason\_t}
\end{itemize}

The default configuration uses \texttt{640x480}, floor \texttt{y = 400},
player size \texttt{48x96}, projectile size \texttt{28x20}, projectile speed
\texttt{6}, walk speed \texttt{3}, jump velocity \texttt{-14}, gravity
\texttt{1}, max HP \texttt{100}, and round duration \texttt{99 * 60} frames.

\subsection{Input Interface}

\texttt{sw/include/fighter\_input.h} defines the input parsing layer.

The default key mapping is:

\begin{itemize}
  \item \texttt{W}: jump
  \item \texttt{A}: move left
  \item \texttt{D}: move right
  \item \texttt{S}: crouch
  \item \texttt{J}: attack
  \item \texttt{K}: guard
  \item \texttt{L}: exit the current match
\end{itemize}

Attack combinations are decoded in software:

\begin{itemize}
  \item \texttt{J}: normal attack
  \item \texttt{A + J}: fireball
  \item \texttt{D + J}: dragon punch
  \item \texttt{W + J}: jump attack
  \item \texttt{W + D + J}: forward jump attack
  \item \texttt{W + A + J}: back jump attack
  \item \texttt{S + J}: sweep
\end{itemize}

\subsection{Rendering Interface}

\texttt{sw/include/fighter\_renderer.h} defines three renderer backends:

\begin{itemize}
  \item \texttt{FIGHTER\_RENDERER\_BACKEND\_MMIO}: FPGA VGA framebuffer IP
  \item \texttt{FIGHTER\_RENDERER\_BACKEND\_FRAMEBUFFER}: Linux framebuffer fallback
  \item \texttt{FIGHTER\_RENDERER\_BACKEND\_CONSOLE}: debug console fallback
\end{itemize}

On Linux, \texttt{fighter\_renderer\_init()} tries the MMIO VGA backend first
when framebuffer-style output is preferred. It then falls back to
\texttt{/dev/fb0}, \texttt{/dev/fb1}, or console output. The renderer loads
menu PPMs, a background PPM, and character sprites from \texttt{game\_assets},
with \texttt{FIGHTER\_ASSET\_ROOT} available as an override.

\subsection{Audio Interface}

\texttt{sw/include/fighter\_audio.h} defines command-based audio playback.

The current tracks are:

\begin{itemize}
  \item \texttt{MENU\_BGM}
  \item \texttt{MENU\_CONFIRM}
  \item \texttt{GAME\_OVER}
\end{itemize}

The supported command types are \texttt{PLAY\_ONCE}, \texttt{START\_LOOP},
and \texttt{STOP\_LOOP}. Audio is disabled by default in the demo. It is
initialized when \texttt{phase1\_demo --audio} or \texttt{input\_demo --audio}
is used. Backend selection is \texttt{WM8731/MMIO} first, then command-line
player fallback, then disabled.

\subsection{USB Keyboard Capture Interface}

\texttt{sw/include/usb\_hid\_keyboard.h} and
\texttt{sw/input/usb\_hid\_keyboard.c} implement device enumeration and polling
through \texttt{libusb}. The code accepts HID boot keyboards, locates the
interrupt IN endpoint, polls 8-byte reports, and supports up to two devices.

\section{Verilog Module Interfaces}

\subsection{\texttt{fighter\_vga\_renderer}}

The module definition is in \texttt{hw/fighter\_vga\_renderer.sv}.

Key interface properties:

\begin{itemize}
  \item clock/reset: \texttt{clk\_50}, \texttt{reset\_n}
  \item Avalon-MM control: \texttt{avs\_chipselect}, \texttt{avs\_read},
  \texttt{avs\_write}
  \item Avalon-MM address/data: \texttt{avs\_address[15:0]},
  \texttt{avs\_writedata[31:0]}, and \texttt{avs\_readdata[31:0]}
  \item VGA conduit: \texttt{vga\_r/g/b[7:0]}, \texttt{vga\_hs},
  \texttt{vga\_vs}, \texttt{vga\_clk}, \texttt{vga\_blank\_n},
  \texttt{vga\_sync\_n}
\end{itemize}

\begin{longtable}{>{\raggedright\arraybackslash}p{0.22\textwidth} >{\raggedright\arraybackslash}p{0.10\textwidth} >{\raggedright\arraybackslash}p{0.12\textwidth} >{\raggedright\arraybackslash}p{0.42\textwidth}}
\toprule
Port & Width & Direction & Meaning \\
\midrule
\endhead
\texttt{clk\_50} & 1 & input & 50 MHz FPGA clock \\
\texttt{reset\_n} & 1 & input & Active-low reset from Platform Designer reset network \\
\texttt{avs\_chipselect} & 1 & input & Avalon slave selected \\
\texttt{avs\_read} & 1 & input & Avalon read strobe \\
\texttt{avs\_write} & 1 & input & Avalon write strobe \\
\texttt{avs\_address} & 16 & input & Word offset for control, geometry, ident, or framebuffer write \\
\texttt{avs\_writedata} & 32 & input & Register write data or packed RGB565 pixels \\
\texttt{avs\_readdata} & 32 & output & Register read data \\
\texttt{vga\_r} & 8 & output & Red channel to VGA DAC \\
\texttt{vga\_g} & 8 & output & Green channel to VGA DAC \\
\texttt{vga\_b} & 8 & output & Blue channel to VGA DAC \\
\texttt{vga\_hs} & 1 & output & VGA horizontal sync \\
\texttt{vga\_vs} & 1 & output & VGA vertical sync \\
\texttt{vga\_clk} & 1 & output & VGA pixel clock \\
\texttt{vga\_blank\_n} & 1 & output & VGA blanking indicator \\
\texttt{vga\_sync\_n} & 1 & output & VGA sync control, fixed low \\
\bottomrule
\end{longtable}

\subsection{\texttt{fighter\_audio\_wm8731}}

The module definition is in \texttt{hw/fighter\_audio.sv}.

Key parameters:

\begin{itemize}
  \item \texttt{FIFO\_DEPTH = 128}
  \item \texttt{CLK\_HZ = 50000000}
  \item \texttt{I2C\_RATE\_HZ = 100000}
  \item \texttt{XCK\_DIV = 4}
  \item \texttt{BCLK\_DIV = 16}
  \item \texttt{I2C\_DEVICE\_ADDR = 7'h1A}
\end{itemize}

Key ports include clock/reset, a 4-word Avalon-MM slave, WM8731 audio pins,
WM8731 I2C pins, and codec initialization status outputs.

\begin{longtable}{>{\raggedright\arraybackslash}p{0.22\textwidth} >{\raggedright\arraybackslash}p{0.10\textwidth} >{\raggedright\arraybackslash}p{0.12\textwidth} >{\raggedright\arraybackslash}p{0.42\textwidth}}
\toprule
Port & Width & Direction & Meaning \\
\midrule
\endhead
\texttt{clk} & 1 & input & 50 MHz FPGA clock \\
\texttt{reset\_n} & 1 & input & Active-low reset \\
\texttt{avs\_chipselect} & 1 & input & Avalon slave selected \\
\texttt{avs\_read} & 1 & input & Avalon read strobe \\
\texttt{avs\_write} & 1 & input & Avalon write strobe \\
\texttt{avs\_address} & 2 & input & Audio register word offset \texttt{0..3} \\
\texttt{avs\_writedata} & 32 & input & Control bits or sample word \\
\texttt{avs\_readdata} & 32 & output & Status or FIFO-space read data \\
\texttt{aud\_xck} & 1 & output & WM8731 master clock \\
\texttt{aud\_bclk} & 1 & output & WM8731 bit clock \\
\texttt{aud\_daclrck} & 1 & output & DAC left/right clock \\
\texttt{aud\_adclrck} & 1 & output & ADC left/right clock \\
\texttt{aud\_dacdat} & 1 & output & DAC serial data \\
\texttt{aud\_adcdat} & 1 & input & ADC serial data, currently unused \\
\texttt{fpga\_i2c\_sclk} & 1 & bidir & Codec I2C clock, open-drain style \\
\texttt{fpga\_i2c\_sdat} & 1 & bidir & Codec I2C data, open-drain style \\
\texttt{codec\_init\_done} & 1 & output & Initialization complete status \\
\texttt{codec\_init\_error} & 1 & output & Initialization failed status \\
\bottomrule
\end{longtable}

\section{Board Resource Use}

The current Quartus build report is
\texttt{hw/output\_files/soc\_system.fit.summary}. It was generated by
Quartus Prime Lite 21.1 for Cyclone V device \texttt{5CSEMA5F31C6}.

\begin{longtable}{>{\raggedright\arraybackslash}p{0.34\textwidth} >{\raggedright\arraybackslash}p{0.20\textwidth} >{\raggedright\arraybackslash}p{0.20\textwidth} >{\raggedright\arraybackslash}p{0.14\textwidth}}
\toprule
Resource & Used & Available & Percent \\
\midrule
\endhead
ALMs & \texttt{1,653} & \texttt{32,070} & \texttt{5\%} \\
Registers & \texttt{1,885} & -- & -- \\
Pins & \texttt{362} & \texttt{457} & \texttt{79\%} \\
Block memory bits & \texttt{8,192} & \texttt{4,065,280} & \texttt{<1\%} \\
RAM blocks / M10Ks & \texttt{2} & \texttt{397} & \texttt{<1\%} \\
DSP blocks & \texttt{2} & \texttt{87} & \texttt{2\%} \\
PLLs & \texttt{0} & \texttt{6} & \texttt{0\%} \\
DLLs & \texttt{1} & \texttt{4} & \texttt{25\%} \\
\bottomrule
\end{longtable}

The fitter hierarchy report attributes the two M10K blocks to
\texttt{fighter\_audio\_wm8731}: one for the left FIFO and one for the right
FIFO. Each FIFO is \texttt{4096} bits, implemented as \texttt{128 x 32}. The
two DSP blocks are attributed to \texttt{fighter\_vga\_renderer} arithmetic
logic. The custom \texttt{fighter\_audio\_wm8731} entity uses approximately
\texttt{164.4} ALMs, \texttt{250} combinational ALUTs, \texttt{332} dedicated
logic registers, \texttt{8192} block memory bits, and \texttt{2} M10Ks. The
custom \texttt{fighter\_vga\_renderer} entity uses approximately
\texttt{1100.3} ALMs, \texttt{1817} combinational ALUTs, \texttt{776}
dedicated logic registers, \texttt{0} M10Ks in the reported fit, and
\texttt{2} DSP blocks.

Board-level resources used by the project include:

\begin{itemize}
  \item HPS DDR3, USB, SD, UART, Ethernet, SPI, and I2C pins through the
  generated HPS subsystem.
  \item VGA output pins: 8-bit red, 8-bit green, 8-bit blue, HS, VS, CLK,
  BLANK\_N, and SYNC\_N.
  \item Audio codec pins: \texttt{AUD\_XCK}, \texttt{AUD\_BCLK},
  \texttt{AUD\_DACLRCK}, \texttt{AUD\_ADCLRCK}, \texttt{AUD\_DACDAT}, and
  \texttt{AUD\_ADCDAT}.
  \item FPGA-side codec I2C pins: \texttt{FPGA\_I2C\_SCLK} and
  \texttt{FPGA\_I2C\_SDAT}.
  \item LEDs: \texttt{LEDR[0]} shows codec initialization done and
  \texttt{LEDR[1]} shows codec initialization error.
\end{itemize}

\section{Implementation Details and Testing}

\subsection{Main Loop and Execution Modes}

\texttt{sw/main\_phase1\_demo.c} is the main demo program. It supports:

\begin{itemize}
  \item \texttt{--script smoke|ko}
  \item \texttt{--usb}
  \item \texttt{--console}
  \item \texttt{--fb PATH}
  \item \texttt{--audio}
  \item \texttt{--frames N}
  \item \texttt{--realtime}
\end{itemize}

At runtime, each logic step collects or synthesizes input reports, parses them
into \texttt{fighter\_player\_result\_t}, updates the game, updates animation
state, and processes audio commands. A rendered frame is drawn after at least
one logic step.

\subsection{Game Rules Already Implemented}

The current branch supports:

\begin{itemize}
  \item left/right movement, jumping, crouching, crouch guard, standing guard
  \item normal attack, fireball, dragon punch, jump attacks, and sweep
  \item fireball projectile spawning, movement, collision, mutual cancellation,
  hit, and block handling
  \item attack startup, active, hit-confirm, block-confirm, and recovery phases
  \item automatic facing updates and grounded overlap resolution
  \item hit damage, guard chip damage, hit stun, block stun, KO, double KO,
  timeout, exit, menu, game over, and victory/KO visual states
\end{itemize}

\subsection{Automated Test Coverage}

\texttt{sw/tests/test\_phase1.c} covers menu entry, match start, exiting to
menu, KO and game-over transitions, restart after game-over gating, timeout
draw resolution, MMIO game-state encoding, hit-confirm behavior, block stun,
airborne cross-up facing reversal, grounded overlap resolution, projectile
behavior, and animation-related state transitions.

\subsection{Build Outputs}

\texttt{sw/Makefile} currently builds:

\begin{itemize}
  \item \texttt{phase1\_demo}
  \item \texttt{phase1\_test}
  \item \texttt{audio\_demo}
  \item \texttt{audio\_probe}
  \item \texttt{vga\_probe}
  \item \texttt{input\_demo}, only when \texttt{libusb-1.0} is installed
\end{itemize}

Recommended validation commands:

\begin{lstlisting}[style=fighterstyle]
cd sw
make
./phase1_test
./vga_probe
./audio_probe
./phase1_demo --script smoke --console
./phase1_demo --script ko --console
./phase1_demo --script smoke --audio
\end{lstlisting}

\subsection{Current Risks and Limitations}

\begin{itemize}
  \item There is no Linux kernel driver yet; VGA and audio hardware are mapped
  directly through \texttt{/dev/mem}.
  \item The current main display path is the FPGA VGA renderer, not Linux
  framebuffer. Its MMIO protocol accepts frame data; it does not yet consume
  the 22-word game-state MMIO encoding directly as object/sprite commands.
  \item Audio is disabled unless requested with \texttt{--audio}.
  \item The code relies on asset files under \texttt{game\_assets}; deployment
  to the board must preserve the expected asset layout or set
  \texttt{FIGHTER\_ASSET\_ROOT}.
  \item \texttt{input\_demo} depends on \texttt{libusb-1.0}; the Makefile only
  builds it when pkg-config can find libusb.
\end{itemize}

\section{Conclusion and Future Work}

The current repository contains a working HPS fighting-game prototype and two
custom FPGA output peripherals: the FPGA VGA renderer and WM8731 audio output.
The important architectural distinction is that the main display path is not
Linux framebuffer; it is a custom FPGA VGA renderer whose current MMIO input
format is frame data rather than object/sprite commands.

The most direct future roadmap is:

\begin{enumerate}
  \item keep the current FPGA VGA renderer path as the stable display backend
  \item add a kernel driver or UIO-style access layer to replace raw
  \texttt{/dev/mem}
  \item optionally add a true FPGA sprite/UI renderer that consumes the
  existing 22-word \texttt{fighter\_mmio\_encode()} contract
  \item expand board-level validation for USB input, VGA, audio, and asset
  loading together
\end{enumerate}

\section{Reference Files}

\begin{itemize}
  \item \texttt{sw/main\_phase1\_demo.c}
  \item \texttt{sw/main\_vga\_probe.c}
  \item \texttt{sw/main\_audio\_probe.c}
  \item \texttt{sw/game/fighter\_game.c}
  \item \texttt{sw/fighter\_animation.c}
  \item \texttt{sw/render\_if/fighter\_renderer.c}
  \item \texttt{sw/render\_if/fighter\_mmio.c}
  \item \texttt{sw/input/fighter\_input.c}
  \item \texttt{sw/input/usb\_hid\_keyboard.c}
  \item \texttt{sw/audio/fighter\_audio.c}
  \item \texttt{sw/include/fighter\_game.h}
  \item \texttt{sw/include/fighter\_renderer.h}
  \item \texttt{sw/include/fighter\_mmio.h}
  \item \texttt{hw/fighter\_vga\_renderer.sv}
  \item \texttt{hw/fighter\_vga\_hw.tcl}
  \item \texttt{hw/fighter\_audio.sv}
  \item \texttt{hw/fighter\_audio\_hw.tcl}
  \item \texttt{hw/soc\_system.qsys}
  \item \texttt{hw/soc\_system\_top.sv}
  \item \texttt{docs/registers.md}
  \item \texttt{docs/hardware\_software\_interface.md}
\end{itemize}

\end{document}
